%# -*- coding:utf-8 -*-
%% start of file `template_en.tex'.
%% Copyright 2006-1008 Xavier Danaux (xdanaux@gmail.com).
%
% This work may be distributed and/or modified under the
% conditions of the LaTeX Project Public License version 1.3c,
% available at http://www.latex-project.org/lppl/.


\documentclass[11pt,a4paper]{moderncv}

\usepackage{fontspec,xunicode}
\usepackage[slantfont,boldfont]{xeCJK}
\usepackage{xcolor}  % replace by the encoding you are using

\usepackage{ifplatform} % 判断平台


% \setmainfont{Tahoma}
\setmainfont{Times New Roman}  % 缺省英文字体.serif是有衬线字体sans serif无衬线字体

\ifwindows
  % Windows 下使用的代码
  \setCJKmainfont[ItalicFont=KaiTi, BoldFont=SimHei]{SimSun}  % 衬线字体 缺省中文字体为
  \setCJKsansfont{SimSun}
  \setmonofont{FangSong}                 % 等宽字体: FangSong, Consolas
\fi

\ifmacosx
  % macOS 下使用的代码
  \setCJKmainfont[ItalicFont={Kai}, BoldFont={Hei}]{STSong}  % 衬线字体 缺省中文字体为
  \setCJKsansfont{STSong}
  \setCJKmonofont{STFangsong}  % 中文等宽字体
\fi


%-----------------------xeCJK下设置中文字体------------------------------%
\setCJKfamilyfont{song}{SimSun}  % 宋体 song
\newcommand{\song}{\CJKfamily{song}}
\setCJKfamilyfont{fs}{FangSong_GB2312}  % 仿宋2312 fs
\newcommand{\fs}{\CJKfamily{fs}}
\setCJKfamilyfont{yh}{Microsoft YaHei}  % 微软雅黑 yh
\newcommand{\yh}{\CJKfamily{yh}}
\setCJKfamilyfont{hei}{SimHei}  % 黑体  hei
\newcommand{\hei}{\CJKfamily{hei}}
\setCJKfamilyfont{hwxh}{STXihei}  % 华文细黑  hwxh
\newcommand{\hwxh}{\CJKfamily{hwxh}}
\setCJKfamilyfont{asong}{Adobe Song Std}  % Adobe 宋体  asong
\newcommand{\asong}{\CJKfamily{asong}}
\setCJKfamilyfont{ahei}{Adobe Heiti Std}  % Adobe 黑体  ahei
\newcommand{\ahei}{\CJKfamily{ahei}}
\setCJKfamilyfont{akai}{Adobe Kaiti Std}  % Adobe 楷体  akai
\newcommand{\akai}{\CJKfamily{akai}}


%------------------------------设置字体大小------------------------%
\newcommand{\chuhao}{\fontsize{42pt}{\baselineskip}\selectfont}  % 初号
\newcommand{\xiaochuhao}{\fontsize{36pt}{\baselineskip}\selectfont}  % 小初号
\newcommand{\yihao}{\fontsize{28pt}{\baselineskip}\selectfont}  % 一号
\newcommand{\erhao}{\fontsize{21pt}{\baselineskip}\selectfont}  % 二号
\newcommand{\xiaoerhao}{\fontsize{18pt}{\baselineskip}\selectfont}  % 小二号
\newcommand{\sanhao}{\fontsize{15.75pt}{\baselineskip}\selectfont}  % 三号
\newcommand{\sihao}{\fontsize{14pt}{\baselineskip}\selectfont}  % 四号
\newcommand{\xiaosihao}{\fontsize{12pt}{\baselineskip}\selectfont}  % 小四号
\newcommand{\wuhao}{\fontsize{10.5pt}{\baselineskip}\selectfont}  % 五号
\newcommand{\subwuhao}{\fontsize{10pt}{\baselineskip}\selectfont}  % 次五号
\newcommand{\xiaowuhao}{\fontsize{9pt}{\baselineskip}\selectfont}  % 小五号
\newcommand{\liuhao}{\fontsize{7.875pt}{\baselineskip}\selectfont}  % 六号
\newcommand{\qihao}{\fontsize{5.25pt}{\baselineskip}\selectfont}  % 七号


% \usepackage{fontawesome}
% \setCJKmainfont[BoldFont={WenQuanYi Micro Hei/Bold}]{WenQuanYi Micro Hei}
% \defaultfontfeatures{Mapping=tex-text}
% \XeTeXlinebreaklocale "zh"
% \XeTeXlinebreakskip = 0pt plus 1pt minus 0.1pt
% moderncv themes
\moderncvtheme[blue]{classic}  % optional argument are 'blue' (default), 'orange', 'red', 'green', 'grey' and 'roman' (for roman fonts, instead of sans serif fonts)
% \moderncvtheme[green]{classic}  % idem
% \moderncvtheme[blue,roman]{hht}
% character encoding



% adjust the page margins
\usepackage[scale=0.9]{geometry}
% \setlength{\hintscolumnwidth}{3cm}  % if you want to change the width of the column with the dates
% \AtBeginDocument{\setlength{\maketitlenamewidth}{6cm}}  % only for the classic theme, if you want to change the width of your name placeholder (to leave more space for your address details
\AtBeginDocument{\recomputelengths}  % required when changes are made to page layout lengths

% personal data
\firstname{毓池}
\familyname{刘}
\title{}  % optional, remove the line if not wanted
% \address{杭州}{}  % optional, remove the line if not wanted
\address{}  % optional, remove the line if not wanted
\mobile{19581520611}  % optional, remove the line if not wanted
% \fax{fax (optional)}  % optional, remove the line if not wanted
\email{834787030@qq.com}  % optional, remove the line if not wanted
\homepage{}  % optional, remove the line if not wanted
\social[github]{}
\extrainfo{%
  求职意向: AI开发工程师/AI安全研究员
}

% \photo[64pt]{avatar.png}  % '64pt' is the height the picture must be resized to and 'picture' is the name of the picture file; optional, remove the line if not wanted
% \quote{China\TeX 您的LaTeX乐园，TeX\&\LaTeX 王国}  % optional, remove the line if not wante

% \nopagenumbers{}  % uncomment to suppress automatic page numbering for CVs longer than one page


%----------------------------------------------------------------------------------
%            content
%----------------------------------------------------------------------------------
\begin{document}
\maketitle
\vspace*{-14mm}

\section{教育经历}
\cventry{2023.09--至今}{硕士}{中南民族大学}{计算机科学与技术}{}{绩点: 4.13/5.0 (5/21)}
\cvlistitem{校级二等奖奖学金 (3次)}
\cvlistitem{泰迪杯国家二等奖}
\cvlistitem{学术成果: MS-ECFormer, IJCNN2026, CCF C, 一作论文 (风电预测, Transformer方向)}
\cventry{2018.09--2022.06}{本科}{中北大学}{计算机科学与技术}{}{绩点: 3.93/5.0 (4/104)}
\cvlistitem{校级二等奖奖学金 (4次)}
\cvlistitem{蓝桥杯国家三等奖}

\section{项目经历}
\cventry{2025.06--2025.09}{智能客服Agent系统}{武汉文弘文创}{大模型应用工程师}{}{%
  针对“步道乐跑”APP数千万大学生用户的咨询需求(如跑量异常、打卡故障等)，为降低人力成本并实现7×24小时精准、安全的服务，独立开发了集成RAG和自动化工具链的智能客服Agent系统，并对Agent的潜在安全风险（如越权工具调用）进行了针对性设计。\newline{}
  项目主要内容:\newline{}
  1. RAG知识库深度优化与安全增强: 针对非结构化《客服操作手册》，采用父子分块策略，父块保留上下文语义，子块细化具体操作步骤。该策略有效缓解了大模型在处理长文本时的“幻觉”问题，将业务知识的检索召回率提升了30\%，并确保回答严格符合业务准则，从源头提升信息输出的安全性与可靠性。\newline{}
  2. 安全导向的Agent架构设计: 利用LangChain框架，结合ReAct推理模式与有限状态机(FSM)。FSM负责锁定并管控涉及用户数据查询、修改等高敏感度业务流程，ReAct负责灵活处理用户咨询意图，通过流程隔离降低未经授权操作的风险。\newline{}
  3. MCP Server工具化封装与安全编排: 将后端数据库查询、打卡记录修复、日志提取等工具封装为标准MCP服务。支持根据用户意图与权限校验结果，通过异步/并发策略动态调用工具，并设计了输入过滤与操作日志记录，以增强系统的可审计性与安全性。
}
\cventry{2024.09--2025.06}{铸牢中华民族共同体知识大模型对话系统}{}{}{}{%
  基于大语言模型构建面向垂直领域的智能对话系统，通过权威数据和前沿技术微调模型，结合RAG和Agent技术，实现多模态交互与复杂任务编排，致力于提供可靠、安全的民族领域智能知识服务。\newline{}
  核心工作:\newline{}
  1. 高质量与高安全性数据集构建: 基于多源异构数据采集和精细化清洗流程，构建了涵盖学术文献、政策文件及官方报道的领域专属语料库。通过分层提示词和自生成与增强生成双路径机制，形成了包含发散性、专业性、权威性三层次的约17万条高质量监督微调数据集，为模型生成内容的准确性与安全性奠定数据基础。\newline{}
  2. 模型微调与可靠性提升: 采用“继续预训练知识注入”和“递进式监督微调”策略。首先利用领域语料库对基座模型进行继续预训练以注入领域知识，随后将微调数据集按层次划分，并基于LoRA技术依次递进微调。最终模型在领域客观任务上较基座模型提升20.4\%，有效缓解了领域知识幻觉与理解偏差，显著提升了生成内容的可靠性与事实准确性。\newline{}
  3. RAG对话模块与事实性防护: 构建了面向垂直领域的自适应双层检索增强生成架构，引入DeepSeek-V3.2作为意图判别路由，采用BGE-M3嵌入模型进行稠密向量化，基于FAISS向量数据库实现关键词聚类分簇存储。通过“关键词匹配限定检索范围+簇内稠密向量相似度计算”的双层检索机制，将Top-5召回准确率提升至87.5\%，MRR达0.734，该设计有效约束了模型生成范围，大幅降低了事实性幻觉风险，是提升AI系统输出安全性的关键实践。
}
\cventry{2024.03--2024.05}{基于视觉语言预训练(VLP)的中文图文检索研究}{}{}{}{%
  针对大规模图文检索中的语义失真和效率瓶颈问题，独立研发了跨模态检索系统。通过两阶段预训练并引入了Flash-Attention和优化存储方案，致力于提升海量异构数据下的检索精度、响应速度及模型鲁棒性。\newline{}
  核心工作:\newline{}
  1. 数据处理、增强与鲁棒性测试: 针对原生数据集图文匹配度低的问题，扩充高质量数据集并利用LMDB存储结构解决小图片读取瓶颈。通过数据增强策略提升模型在复杂、噪声场景下的鲁棒性，并对模型处理对抗性样本的稳定性进行了初步测试。\newline{}
  2. 训练策略与性能优化: 采用两阶段训练逻辑，结合Random Masking和Flash-Attention技术进行算子级优化，提升训练效率。\newline{}
  3. 模型性能与鲁棒性调优: 以Recall@1/5/10为核心优化目标，通过细粒度图像文本特征对齐，使模型在处理包含位置关系和否定信息等复杂语义时的准确率提升了12.39\%，增强了模型对非常规输入的语义理解能力。
}

\section{专业技能}
\cvline{\textbf{编程与开发}}{掌握Python/C++，具备扎实的数据结构与算法基础；具备大模型应用开发与RAG优化经验，熟悉LangChain、LlamaIndex等主流智能体框架及其在构建安全、可控Agent流程中的应用。}
\cvline{\textbf{大模型与AI安全}}{了解ReAct、Reflection等Agent推理范式以及AutoGen多智能体框架的基本原理；具备大模型微调（LoRA、SFT）及继续预训练实战经验，关注大模型安全与对齐领域，对提示安全、数据隐私保护等技术有初步了解。}
\cvline{\textbf{工具与综合能力}}{熟练运用Git进行分支管理与团队协作；具备出色的技术问题拆解能力，能快速跟进AI安全等前沿领域的技术动态，学习与适应能力强。}

\section{学术成果}
\cvline{已录用}{MS-ECFormer, IJCNN2026, CCF C, 一作论文 (风电预测, Transformer方向)}

\end{document}


%% end of file `template_en.tex'.

%%% Local Variables:
%%% mode: latex
%%% TeX-command-extra-options: "-shell-escape"
%%% TeX-master: t
%%% TeX-engine: xetex
%%% End: